\section{Conclusion}
\label{sec:conclusion}

This project evaluated whether a Decision Transformer can provide useful warm-start trajectories for nonlinear minimum-time vehicle trajectory optimization with obstacle avoidance. The resulting system combined offline sequence modeling of optimizer-generated trajectories with autoregressive rollout, wrapper-based validation and projection, and downstream nonlinear solver refinement.

Empirically, larger-context training and targeted repair-data fine-tuning substantially improved offline validation performance. However, these improvements did not produce a convincing solver-facing advantage in the corrected deployment benchmarks. Under the default gate, acceptance remained near zero; under relaxed acceptance, some trajectories were admitted, but they still did not reduce solver work on average. The final evidence therefore indicates that offline sequence-model quality was not a reliable proxy for warm-start utility in this setting.

The most plausible remaining bottleneck is early lateral rollout behavior. The learned model often failed not because of global incoherence, but because it entered the wrong local region of the trajectory space too early for the downstream solver to benefit. Future work should therefore focus less on improving average offline prediction quality and more on directly targeting deployment-critical early-horizon decisions and solver-relevant rollout geometry.